
\section{时间有序积模空间与逐时积模空间的定向}
\label{sec:8-9}
\liftIndex{I307}

\subsection{\texorpdfstring{$\mathcal M_{k+1}^{\mathrm{main}}
(M',L',\{J_\rho\}_\rho:\beta;\operatorname{top}(\rho))$ 的定向}%
{top 模空间的定向}}
\label{subsec:8-9-1}

在第~\ref{sec:4-6} 节的定理~\ref{thm:4-6-1} 中\translatornote{源书
此处仍将正式编号为定理的 4.6.1 误称为命题；依其正式陈述改正。}，
我们用模空间
\[
 \mathcal M_{k+1}^{\mathrm{main}}
 (M',L',\{J_\rho\}_\rho:\beta;\operatorname{top}(\rho))
\]
定义滤过 $A_\infty$ 同态。一般情形可约化到辛微分同胚
$\varphi$ 为恒等映射，而 $\{J_\rho\}_\rho$ 是任意一族与
$\omega$ 相容的殆复结构的情形。（不过，为了分清定义域与靶，我们仍使用
$(M,L)$ 和 $(M',L')$。）本小节的目标是为
\[
\begin{aligned}
 &\mathcal M_{k+1}^{\mathrm{main}}
 (\{J_\rho\}_\rho:\beta;\operatorname{top}(\rho);P_1,\ldots,P_k)\\
 &\quad:=
 \mathcal M_{k+1}^{\mathrm{main}}
 (M',L',\{J_\rho\}_\rho:\beta;
  \operatorname{top}(\rho);P_1,\ldots,P_k)
\end{aligned}
\]
赋予典范定向，使所得 $\mathfrak f$ 成为滤过 $A_\infty$ 同态。

先回忆
\[
 \mathcal M_{k+1}^{\mathrm{main}}
 (M',L',\{J_\rho\}_\rho:\beta;\operatorname{top}(\rho))
\]
的一个元素。设 $((\Sigma,\vec z),(u_\alpha),(\rho_\alpha))$
是满足下列条件的系统：
\begin{enumerate}[label=\textup{(4.6.7.\arabic*)},leftmargin=5.9em]
\item
$u_\alpha\colon(\Sigma_\alpha,\partial\Sigma_\alpha)\to(M,L)$
是 $J_{\rho_\alpha}$-全纯映射。
\item
$\rho_\alpha\in[0,1]$；若 $\alpha_1\leq\alpha_2$，则
$\rho_{\alpha_1}\leq\rho_{\alpha_2}$。
\item
$((\Sigma,\vec z),(u_\alpha))$ 在定义~\ref{def:2-1-24}
的意义下是稳定的。
\item
$u_\alpha$ 的同调类为 $\beta(\alpha)$，且
$\sum_\alpha\beta(\alpha)=\beta$。
\item
若 $z\in\Sigma_\alpha\cap\Sigma_{\alpha'}$，则
$u_\alpha(z)=u_{\alpha'}(z)$。
\end{enumerate}

以
\[
 \mathcal M_{k+1}^{\mathrm{main}}
 (\beta;P_1,\ldots,P_k;J_\rho)
\]
表示相对于殆复结构 $J_\rho$、由式~(3.5.1) 中对象构成的模空间。
模空间
\[
 \mathcal M_{k+1}^{\mathrm{main}}
 (M',L',\{J_\rho\}_\rho:\beta;\operatorname{top}(\rho))
\]
的一个顶维层由式~(4.6.7) 中定义域不可约的对象构成。定义域不可约时，
把它定向为
\begin{equation}
 (-1)^{n+1}
 \coprod_{0\leq\rho\leq1}
 \{\rho\}\times
 \mathcal M_{k+1}^{\mathrm{main},\mathrm{reg}}
 (\beta;P_1,\ldots,P_k;J_\rho),
\tag{8.9.1}\label{eq:8-9-1}
\end{equation}
并将其记为
\[
\begin{aligned}
 &\mathcal M_{k+1}^{\mathrm{main},\mathrm{reg}}
 (\{J_\rho\}_\rho:\beta;\operatorname{top}(\rho);
  P_1,\ldots,P_k)\\
 &\qquad\subset
 \mathcal M_{k+1}^{\mathrm{main}}
 (\{J_\rho\}_\rho:\beta;\operatorname{top}(\rho);
  P_1,\ldots,P_k).
\end{aligned}
\]
回忆，对每个固定的 $J_\rho$，
\[
\begin{aligned}
 &\mathcal M_{k+1}^{\mathrm{main},\mathrm{reg}}
 (\beta;P_1,\ldots,P_k;J_\rho)\\
 &\quad={}
 (-1)^{\epsilon_1}
 \mathcal M_{k+1}^{\mathrm{main},\mathrm{reg}}
 (L;\beta;J_\rho)_{
 (\operatorname{ev}_1,\ldots,\operatorname{ev}_k)}
 \mathop{\times}_{f_1\times\cdots\times f_k}
 \left(\prod_{i=1}^{k}P_i\right),
\end{aligned}
\]
其中
\[
 \epsilon_1=(n+1)
 \sum_{j=1}^{k-1}\sum_{i=1}^{j}\deg P_i.
\]
见定义~\ref{def:8-4-1}。注意，出现在
$\mathcal M_{k+1}^{\mathrm{main},\mathrm{reg}}
(\beta;P_1,\ldots,P_k;J_\rho)$ 中的 $P_i$，
通过 $f_i\colon P_i\to L$ 被看作 $L$ 中的奇异单形。

命题~\ref{prop:8-5-1}(1) 蕴含
\[
\begin{aligned}
 &(-1)^{\epsilon'_2+1}
 \coprod_{0\leq\rho\leq1}\{\rho\}\times
 \mathcal M_{k-\ell+2}^{\mathrm{main},\mathrm{reg}}
 \Bigl(\beta';P_1,\ldots,P_{i-1},\\[-2pt]
 &\hspace{8em}
 \mathcal M_{\ell+1}^{\mathrm{main},\mathrm{reg}}
 (\beta'';P_i,\ldots,P_{i+\ell-1};J_\rho),
 P_{i+\ell},\ldots,P_k;J_\rho\Bigr)\\
 &\quad\subset
 \partial\mathcal M_{k+1}^{\mathrm{main},\mathrm{reg}}
 (\{J_\rho\}_\rho:\beta;\operatorname{top}(\rho);
 P_1,\ldots,P_k),
\end{aligned}
\]
其中
\[
 \epsilon'_2=\epsilon_2-(n+1)
 =\sum_{j=1}^{i-1}(\deg P_j+1),
 \qquad \beta=\beta'+\beta''.
\]
另一方面，命题~\ref{prop:8-5-1}(2) 蕴含
\[
\begin{aligned}
 &\coprod_{0\leq\rho\leq1}\{\rho\}\times
 \mathcal M_{k-\ell+2}^{\mathrm{main},\mathrm{reg}}
 \Bigl(\beta';P_1,\ldots,P_{i-1},\\[-2pt]
 &\hspace{8em}
 \mathcal M_{\ell+1}^{\mathrm{main},\mathrm{reg}}
 (\beta'';P_i,\ldots,P_{i+\ell-1};J_\rho),
 P_{i+\ell},\ldots,P_k;J_\rho\Bigr)\\
 &\quad\subset
 \mathcal M_{k-\ell+2}^{\mathrm{main},\mathrm{reg}}
 \Bigl(\{J_\rho\}_\rho:\beta';\operatorname{top}(\rho);
 P_1,\ldots,P_{i-1},\\[-2pt]
 &\hspace{12em}
 \partial Q_{i,\ell}^{\rho}(\beta''),
 P_{i+\ell},\ldots,P_k\Bigr)\\
 &\quad\subset
 (-1)^{\epsilon_3+1}
 \partial\mathcal M_{k-\ell+2}^{\mathrm{main}}
 \Bigl(\{J_\rho\}_\rho:\beta';\operatorname{top}(\rho);
 P_1,\ldots,P_{i-1},\\[-2pt]
 &\hspace{12em}
 Q_{i,\ell}^{\rho}(\beta''),
 P_{i+\ell},\ldots,P_k\Bigr),
\end{aligned}
\]
其中\translatornote{源书在上式后两个外层模空间中均把类写成
$\beta$；但第一项的外层类是 $\beta'$、内层类是 $\beta''$，且紧前
已规定 $\beta=\beta'+\beta''$。依命题 8.5.1(1) 的同一边界分解，
两处均改作 $\beta'$。}
\[
 \epsilon_3=1+\sum_{j=1}^{i-1}(\deg P_j+1),
\]
并且
\begin{equation}
 Q_{i,\ell}^{\rho}(\beta'')
 :=(-1)^{n+1}
 \coprod_{0\leq\sigma\leq\rho}
 \{\sigma\}\times
 \mathcal M_{\ell+1}^{\mathrm{main},\mathrm{reg}}
 (\beta'';P_i,\ldots,P_{i+\ell-1};J_\sigma).
\tag{8.9.2}\label{eq:8-9-2}
\end{equation}

因此，
\begin{equation}
\begin{aligned}
 &\coprod_{0\leq\rho\leq1}\{\rho\}\times
 \mathcal M_{k-\ell+2}^{\mathrm{main},\mathrm{reg}}
 \Bigl(\beta';P_1,\ldots,P_{i-1},\\[-2pt]
 &\hspace{8em}
 \mathcal M_{\ell+1}^{\mathrm{main},\mathrm{reg}}
 (\beta'';P_i,\ldots,P_{i+\ell-1};J_\rho),
 P_{i+\ell},\ldots,P_k;J_\rho\Bigr)
\end{aligned}
\tag{8.9.3}\label{eq:8-9-3}
\end{equation}
作为
\[
 \mathcal M_{k+1}^{\mathrm{main},\mathrm{reg}}
 (\{J_\rho\}_\rho:\beta;\operatorname{top}(\rho);
 P_1,\ldots,P_k)
\]
的边界时所得的定向，与它作为
\[
\begin{aligned}
 \mathcal M_{k-\ell+2}^{\mathrm{main}}
 \bigl(&\{J_\rho\}_\rho:\beta';\operatorname{top}(\rho);
 P_1,\ldots,P_{i-1},\\[-2pt]
 &Q_{i,\ell}^{\rho}(\beta''),
 P_{i+\ell},\ldots,P_k\bigr)
\end{aligned}
\]
的边界时所得的定向相反。因此，
\[
 \mathcal M_{k+1}^{\mathrm{main},\mathrm{reg}}
 (\{J_\rho\}_\rho:\beta;\operatorname{top}(\rho);
 P_1,\ldots,P_k)
\]
上的定向与
\[
\begin{aligned}
 \mathcal M_{k-\ell+2}^{\mathrm{main}}
 \bigl(&\{J_\rho\}_\rho:\beta';\operatorname{top}(\rho);
 P_1,\ldots,P_{i-1},\\[-2pt]
 &Q_{i,\ell}^{\rho}(\beta''),
 P_{i+\ell},\ldots,P_k\bigr)
\end{aligned}
\]
上的定向，在余维 $1$ 的层~\eqref{eq:8-9-3} 上相合。

更一般地，对定义~\ref{def:4-6-8} 中对象构成的模空间，我们按奇点数，
也就是不可约分支数，归纳定义
\[
 \mathcal M_{k+1}^{\mathrm{main}}
 (\{J_\rho\}_\rho:\beta;\operatorname{top}(\rho);
 P_1,\ldots,P_k)
\]
上的定向。对 $c\in[0,1]$ 和 $(\beta,k)\ne(\beta_0,1)$，以
\[
 \mathcal M_{k+1}^{\mathrm{main},\leq c}
 (\{J_\rho\}_\rho:\beta;\operatorname{top}(\rho);
 P_1,\ldots,P_k)
\]
表示式~(4.6.7) 中所有 $\rho_\alpha\in[0,c]$ 的对象所构成的模空间。

假设对一切定义域至多含 $d$ 个不可约分支的元素以及一切
$c\in[0,1]$，上述模空间的定向已经确定。令 $u$ 是式~(4.6.7)
中带 $d+1$ 个不可约分支的对象。于是 $u$ 属于
\[
 \{\sigma\}\times
 \mathcal M_{\ell+1}^{\mathrm{main},\mathrm{reg}}
 (\beta';Q_1,\ldots,Q_\ell;J_\sigma),
\]
其中每个 $Q_j$ 或者是 (a) 某个 $P_i$，或者是 (b) 下列空间中的
元素\translatornote{源书在此使用下标 $k_{i_j}+1$，但本段没有定义
$k_{i_j}$，且显示的输入块由 $i_{j-1}+1,\ldots,i_j$ 编号；现有上下文
不足以唯一判断这是漏定义还是下标误排，故暂按源式保留。}：
\[
 \mathcal M_{k_{i_j}+1}^{\mathrm{main},\leq\sigma}
 \bigl(\{J_\rho\}_\rho:\beta_{i_j};\operatorname{top}(\rho);
 P_{i_{j-1}+1},\ldots,P_{i_j}\bigr).
\]
由假设，至少有一个 $Q_j$ 属于 (b) 型。还要注意，(b) 型 $Q_j$
所代表元素中的不可约分支数不超过 $d$。因而归纳假设已为所有 $Q_j$
赋予定向。于是，对
\[
 \coprod_{0\leq\rho\leq c}\{\rho\}\times
 \mathcal M_{\ell+1}^{\mathrm{main},\mathrm{reg}}
 (\beta';Q_1,\ldots,Q_\ell;J_\rho)
\]
采用式~\eqref{eq:8-9-1} 的定向；这就在 $u$ 处确定了
\[
 \mathcal M_{k+1}^{\mathrm{main},\leq c}
 (\{J_\rho\}_\rho:\beta;\operatorname{top}(\rho);
 P_1,\ldots,P_k)
\]
的定向。

同上述讨论（$d=2$ 的情形）一样，可以看出这些定向在余维 $1$ 的层上
相合，从而为模空间
\[
 \mathcal M_{k+1}^{\mathrm{main}}
 (M',L',\{J_\rho\}_\rho:\beta;\operatorname{top}(\rho);
 P_1,\ldots,P_k)
\]
给出相容定向。

下面讨论定理~\ref{thm:4-6-1} 与命题~\ref{prop:4-6-14} 中映射
$\widehat{\mathfrak f}$ 的符号问题。先陈述如下结果。

\begin{proposition}{8.9.4}\label{prop:8-9-4}
设 $(\beta,k)\ne(\beta_0,1)$。模空间
\[
 \mathcal M_{k+1}^{\mathrm{main}}
 (\{J_\rho\}_\rho:\beta;\operatorname{top}(\rho);
 P_1,\ldots,P_k)
\]
的边界是下列各层的闭包。令 $k_\alpha+1$ 为 $\Sigma_\alpha$ 上特殊点
（即标记点或奇点）的个数。

\par\smallskip\noindent
\textup{情形 (1).}
$\rho_{\alpha_0}=1$ 且 $(\beta',\ell)\ne(\beta_0,1)$。有
\[
\begin{aligned}
 &(-1)^{n+1}
 \mathcal M_{\ell+1}^{\mathrm{main}}
 (\beta';Q_1,\ldots,Q_\ell;J')\\
 &\quad\subset
 \partial\mathcal M_{k+1}^{\mathrm{main}}
 (\{J_\rho\}_\rho:\beta;\operatorname{top}(\rho);
 P_1,\ldots,P_k),
\end{aligned}
\]
其中 $Q_j$ 或者是某个 $P_i$（此时记 $\beta^{(j)}=\beta_0$），或者是
下列空间\translatornote{源书把此输入块的末项排成 $P_j$；同一构造
在上一页写作 $P_{i_{j-1}+1},\ldots,P_{i_j}$，而式 (3.2.28) 也按连续
输入块定义 $\widehat{\mathfrak f}$，故改为 $P_{i_j}$。}：
\[
 Q_j=
 \mathcal M_{k_j+1}^{\mathrm{main}}
 \bigl(\{J_\rho\}_\rho:\beta^{(j)};\operatorname{top}(\rho);
 P_{i_{j-1}+1},\ldots,P_{i_j}\bigr),
\]
这里 $\beta^{(j)}\ne\beta_0$，且
\[
 \beta=\beta'+\sum_j\beta^{(j)}.
\]

\par\smallskip\noindent
\textup{情形 (2).}
某个 $\rho_\alpha=0$。有
\[
\begin{aligned}
 &(-1)^{n+1+\epsilon_3}
 \mathcal M_{k-\ell+2}^{\mathrm{main}}
 \Bigl(\{J_\rho\}_\rho:\beta';\operatorname{top}(\rho);
 P_1,\ldots,P_{i-1},\\[-2pt]
 &\hspace{9em}
 \mathcal M_{\ell+1}^{\mathrm{main}}
 (\beta'';P_i,\ldots,P_{i+\ell-1}),
 P_{i+\ell},\ldots,P_k\Bigr)\\
 &\quad\subset
 \partial\mathcal M_{k+1}^{\mathrm{main}}
 (\{J_\rho\}_\rho:\beta;\operatorname{top}(\rho);
 P_1,\ldots,P_k),
\end{aligned}
\]
其中
\[
 \beta''=\sum_{\alpha'\leq\alpha}\beta(\alpha'),
 \qquad
 \epsilon_3=1+\sum_{j=1}^{i-1}(\deg P_j+1).
\]

\par\smallskip\noindent
\textup{情形 (3).}
把 $P_i$ 换成其边界 $\partial P_i$。有
\[
\begin{aligned}
 &(-1)^{1+\epsilon_3}
 \mathcal M_{k+1}^{\mathrm{main}}
 (\{J_\rho\}_\rho:\beta;\operatorname{top}(\rho);
 P_1,\ldots,P_{i-1},\partial P_i,P_{i+1},\ldots,P_k)\\
 &\quad\subset
 \partial\mathcal M_{k+1}^{\mathrm{main}}
 (\{J_\rho\}_\rho:\beta;\operatorname{top}(\rho);
 P_1,\ldots,P_k).
\end{aligned}
\]

\par\smallskip\noindent
\textup{情形 (4).}
$\rho_{\alpha_0}=0$。有
\[
\begin{aligned}
 &(-1)^n
 \mathcal M_{k+1}^{\mathrm{main}}
 (\beta;P_1,\ldots,P_k;J_0)\\
 &\quad\subset
 \partial\mathcal M_{k+1}^{\mathrm{main}}
 (\{J_\rho\}_\rho:\beta;\operatorname{top}(\rho);
 P_1,\ldots,P_k).
\end{aligned}
\]
\end{proposition}

\begin{proof}
在情形 (2) 中，由稳定性条件，
\[
 (\beta(\alpha),k_\alpha)\ne(\beta_0,1),
 \qquad
 (\beta',k-\ell+1)\ne(\beta_0,1).
\]
情形 (3) 与 (4) 可以分别看成
\[
 (\beta(\alpha),k_\alpha)=(\beta_0,1),
 \qquad
 (\beta',k-\ell+1)=(\beta_0,1)
\]
的情形；也就是说，$L$ 与 $L'$ 上的经典边界算子
$\partial=(-1)^n\mathfrak m_{1,0}$ 分别出现在情形 (3) 与 (4) 中。
证明本命题时，只须注意
\[
\begin{aligned}
 &(-1)^{n+1}
 \mathcal M_{k+1}^{\mathrm{main}}
 (\beta;P_1,\ldots,P_k;J_1)\\
 &\quad\subset
 \partial\mathcal M_{k+1}^{\mathrm{main}}
 (\{J_\rho\}_\rho:\beta;\operatorname{top}(\rho);
 P_1,\ldots,P_k),\\[4pt]
 &(-1)^n
 \mathcal M_{k+1}^{\mathrm{main}}
 (\beta;P_1,\ldots,P_k;J_0)\\
 &\quad\subset
 \partial\mathcal M_{k+1}^{\mathrm{main}}
 (\{J_\rho\}_\rho:\beta;\operatorname{top}(\rho);
 P_1,\ldots,P_k);
\end{aligned}
\]
见式~\eqref{eq:8-9-1} 与约定~\ref{conv:8-2-1}(1)。随后便可用与命题
\ref{prop:8-5-1} 类似的方法证明命题~\ref{prop:8-9-4}。
\end{proof}

对 $(k,\beta)\ne(1,\beta_0)$，置
\[
\begin{aligned}
 \mathfrak f_{k,\beta}(P_1,\ldots,P_k)
 =\Bigl(&\mathcal M_{k+1}^{\mathrm{main}}
 (M',L',\{J_\rho\}_\rho:\beta;\operatorname{top}(\rho);
 P_1,\ldots,P_k),\operatorname{ev}_0\Bigr)
\end{aligned}
\]
并置
\[
 \mathfrak f_{1,\beta_0}(P)=P.
\]
显然，
\[
 \mathfrak f_{1,\beta_0}\circ\mathfrak m_{1,\beta_0}
 =\mathfrak m_{1,\beta_0}\circ\mathfrak f_{1,\beta_0}.
\]
情形 (1) 的各层对应于
\[
 (-1)^{n+1}
 (\mathfrak m_{\ell,\beta'}\circ
  \widehat{\mathfrak f}_{\beta''})(P_1,\ldots,P_k),
\]
其中 $\beta=\beta'+\beta''$ 且
$(\beta',\ell)\ne(\beta_0,1)$。这里回忆
\[
 \widehat{\mathfrak f}_{\beta''}
 =\sum
 \mathfrak f_{i_1,\beta^{(1)}}\otimes\cdots\otimes
 \mathfrak f_{i_\ell,\beta^{(\ell)}},
 \qquad
 \beta''=\beta^{(1)}+\cdots+\beta^{(\ell)};
\]
见式~\eqref{eq:3-2-28}。情形 (2) 的各层对应于
\[
 (-1)^n
 (\mathfrak f_{k-\ell+1,\beta'}\circ
  \widehat{\mathfrak m}_{\ell,\beta''})(P_1,\ldots,P_k),
\]
其中
\[
 (\beta',k-\ell+1)\ne(\beta_0,1),
 \qquad
 (\beta'',\ell)\ne(\beta_0,1).
\]
这里 $\widehat{\mathfrak m}_{\ell,\beta''}$ 是把
$\mathfrak m_{\ell,\beta''}$ 扩张成分次余导子所得的算子；见式
\eqref{eq:3-2-15}。情形 (3) 的各层对应于
\[
 (-1)^n
 (\mathfrak f_{k,\beta}\circ
  \widehat{\mathfrak m}_{1,\beta_0})(P_1,\ldots,P_k),
\]
而情形 (4) 的各层对应于
\[
 (-1)^n
 (\mathfrak f_{1,\beta_0}\circ
  \mathfrak m_{k,\beta})(P_1,\ldots,P_k).
\]
另一方面，
\[
 \partial\mathcal M_{k+1}^{\mathrm{main}}
 (\{J_\rho\}_\rho:\beta;\operatorname{top}(\rho);
 P_1,\ldots,P_k)
\]
对应于
\[
 (-1)^n
 (\mathfrak m_{1,\beta_0}\circ
  \mathfrak f_{k,\beta})(P_1,\ldots,P_k).
\]
因而，计入
$\mathfrak f_{1,\beta_0}\circ\mathfrak m_{1,\beta_0}
=\mathfrak m_{1,\beta_0}\circ\mathfrak f_{1,\beta_0}$ 后，得到
\[
 \mathfrak m\circ\widehat{\mathfrak f}
 =\widehat{\mathfrak f}\circ\widehat d.
\]

最后，为前面各章用到的含逐时积的模空间定义定向。对模空间
\[
\begin{aligned}
 &\mathcal M_{k+1}^{\mathrm{main}}
 (M',L',\{J_{1,s}\}_s;\beta;\operatorname{twp}(s))\\
 &\quad=
 \coprod_{s\in[0,1]}\{s\}\times
 \mathcal M_{k+1}^{\mathrm{main}}
 ((L',J_{1,s}),\beta),
\end{aligned}
\]
把参数 $s\in[0,1]$ 放在模空间
$\mathcal M_{k+1}^{\mathrm{main}}((L',J_{1,s}),\beta)$ 之前，
由此赋予定向。类似地，对模空间
\[
\begin{aligned}
 &\mathcal M_{k+1}^{\mathrm{main}}
 (M',L',\{J_{\rho,s}\}_{\rho,s};\beta;
  \operatorname{top}(\rho),\operatorname{twp}(s))\\
 &\quad=
 \coprod_{s\in[0,1]}\{s\}\times
 \mathcal M_{k+1}^{\mathrm{main}}
 (M',L',\{J_{\rho,s}\}_\rho:\beta;
  \operatorname{top}(\rho)),
\end{aligned}
\]
把参数 $s\in[0,1]$ 放在模空间
\[
 \mathcal M_{k+1}^{\mathrm{main}}
 (M',L',\{J_{\rho,s}\}_\rho:\beta;\operatorname{top}(\rho))
\]
之前，由此赋予定向。
